%----------------------------------------------------------------------------
%-----UOL_Praesentation------------------------------------------------------
%-----Erstellt von Natalia Chernov im Auftrag der CvO Universität Oldenburg--
%----------------------------------------------------------------------------
%--Diese Vorlage basiert auf der Klasse 'beamer', für weitere Informationen-- 
%--siehe--------------------------------------------------------------------- 
%-http://mirrors.ctan.org/macros/latex/contrib/beamer/doc/beameruserguide.pdf
%----------------------------------------------------------------------------
%--Kompilieren mit 'pdflatex UOL_Praesentation.tex'--------------------------
%--Die Datei 'UOL_Logo_Negativ.png' soll sich im selben Verzeichnis befinden-
%----------------------------------------------------------------------------
\documentclass[xcolor=table,9pt,aspectratio=169]{beamer}

\usepackage[english, ngerman]{babel} 
\usepackage{lmodern}
\usepackage[utf8]{inputenc} 
\usepackage{helvet}
\renewcommand{\familydefault}{\sfdefault}
\usepackage{tikz}
\usepackage{datetime}
\usepackage{anyfontsize}
\usepackage{lipsum}
\usepackage{multicol} 
\usepackage{smartdiagram}

\definecolor{uolblue}{RGB}{0,62,107}

\definecolor{blue1}{RGB}{0,78,159} 
\definecolor{blue2}{RGB}{0,171,217}
\definecolor{blue3}{RGB}{91,197,242}
\definecolor{blue4}{RGB}{161,217,248}

\definecolor{green1}{RGB}{0,120,120}
\definecolor{green2}{RGB}{0,168,121}
\definecolor{green3}{RGB}{148,193,28}
\definecolor{green4}{RGB}{199,211,0}

\definecolor{orange1}{RGB}{213,59,10}
\definecolor{orange2}{RGB}{238,113,0}
\definecolor{orange3}{RGB}{243,145,0}
\definecolor{orange4}{RGB}{253,195,0}

\definecolor{gr}{RGB}{191,191,191}

\setbeamertemplate{frametitle}{\color{uolblue}\fontsize{12}{20}\selectfont{\insertframetitle}}

\pgfdeclareimage[width=0.145\paperwidth]{UOLLogo}{UOL_Logo_Negativ}
\pgfdeclareimage[width=0.072\paperwidth]{UOLLogo_small}{UOL_Logo_Negativ}

\defbeamertemplate*{background canvas}{default_page}
{%
\begin{tikzpicture}
	\useasboundingbox (0,0) rectangle (\the\paperwidth, \the\paperheight);        
	\filldraw[fill=uolblue,fill opacity=1,draw=none] (0,0) rectangle (0.119\paperwidth, \the\paperheight);
    \filldraw[fill=blue2,fill opacity=1,draw=none] (0.119\paperwidth, 0) -- (0.119\paperwidth, 0.565\paperheight) arc (117.2:180:0.6\paperwidth) -- cycle;
    \pgftext[at=\pgfpoint{10}{\the\paperheight-11.5}, left, top]{\pgfsetfillopacity{1}\pgfuseimage{UOLLogo_small}}; 
\end{tikzpicture}
}
\defbeamertemplate*{background canvas}{titlepage_image}
{
\begin{tikzpicture}
    \useasboundingbox (0,0) rectangle (\the\paperwidth, \the\paperheight); 
	\filldraw[fill=uolblue,fill opacity=1, draw=none] (0,0) rectangle (\the\paperwidth, \the\paperheight);
    \filldraw[fill=blue2, fill opacity=1, draw=none] (\the\paperwidth, 0) -- (\the\paperwidth, 0.66\paperheight) arc (90:180:0.6\paperwidth) -- cycle;
    \pgftext[at=\pgfpoint{14}{\the\paperheight-17.5}, left, top]{\pgfsetfillopacity{1}\pgfuseimage{UOLLogo}}; 
\end{tikzpicture}  
}
\BeforeBeginEnvironment{frame}{%
  \setbeamertemplate{background canvas}[default_page]%
}
\makeatletter
\define@key{beamerframe}{titlepage_image}[true]{%
  \setbeamercovered{invisible}%
  \setbeamertemplate{background canvas}[titlepage_image]%
 }
\makeatother%

 \setbeamertemplate{footline}
{
  \leavevmode
  \hbox{
  \hspace*{.025\paperwidth}\begin{beamercolorbox}[wd=.094\paperwidth,ht=2.25ex,dp=1ex,left]{}
  ~
  
  \vspace*{.042\paperheight}
     \fontsize{4.4}{5.9} \selectfont \color{white} \textbf{Seite \insertframenumber} \newline \insertdate
  \vspace*{.026\paperheight}
  \end{beamercolorbox}
  \hspace*{.05\paperwidth}\begin{beamercolorbox}
  [wd=.79\paperwidth,ht=2.25ex,dp=1ex,left]{}  
 ~ 

  \vspace*{.042\paperheight}
     \fontsize{4.4}{5.9} \selectfont \color{black} \textbf{\inserttitle}~---~\insertsubtitle \newline \color{gray}Titel \insertauthor~---~Einrichtung
  \vspace*{.026\paperheight}
  \end{beamercolorbox} 
  }
  \vskip0pt
}
%--------------------------------------------------------------------------------------------
\setbeamerfont{title}{size={\fontsize{22}{25}}}
\setbeamerfont{subtitle}{size={\fontsize{12}{14}}}
\setbeamerfont{author}{size={\fontsize{9}{11}}}
\setbeamerfont{date}{size={\fontsize{9}{11}}}
\setbeamercolor{title}{fg=white} 
\setbeamercolor{subtitle}{fg=white} 
\setbeamercolor{author}{fg=white} 
\setbeamercolor{date}{fg=white} 
\setbeamercolor{color_Logo-Platzhalter}{fg=white,bg=gray!40}

\defbeamertemplate*{title page}{customized}[1][]
{ \vspace*{20mm}
  \hspace*{-22.5mm}
  \begin{minipage}{\textwidth}
  \usebeamerfont{title}\usebeamercolor[fg]{title}\inserttitle\par
  \bigskip
  \usebeamerfont{subtitle}\usebeamercolor[fg]{subtitle}\insertsubtitle\par
  \bigskip
  \usebeamerfont{author}\usebeamercolor[fg]{author}\insertauthor, 
  \usebeamerfont{date}\usebeamercolor[fg]{date}\insertdate\par
  \end{minipage}
}
\setbeamertemplate{navigation symbols}{}
\setbeamersize{text margin left=0.17\paperwidth,text margin right=0.04\paperwidth} 

\title{Titel der Präsentation}
\subtitle{Untertitel der Präsentation}
\author{Vorname Name}
\date{\renewcommand{\dateseparator}{.} \ddmmyyyydate \today}
\usepackage{enumitem}
\def\labelitemi{--}
\def\labelitemii{--}
\def\labelitemiii{--}

\begin{document}
%---------------------------------------------------------------------------------
%-----Titelseite------------------------------------------------------------------
%--------------------------------------------------------------------------------- 
{    
\setbeamertemplate{footline}{} 
  \begin{frame}[titlepage_image]
  \maketitle
   
  \end{frame}
}
%---------------------------------------------------------------------------------
%-----Seite 2---------------------------------------------------------------------
%---------------------------------------------------------------------------------  
\begin{frame}
 \begin{overlayarea}{\textwidth}{0.81\paperheight}
{

\vspace*{11mm}

\usebeamerfont{title} \textcolor{uolblue}
{Überschrift Kapiteltrenner}}

\vspace*{7mm}

{\usebeamerfont{subtitle} \textcolor{uolblue}
{Untertitel Kapiteltrenner}}
 
\end{overlayarea}
\begin{overlayarea}{\textwidth}{0.1\paperheight}

    \begin{minipage}{0.5\textwidth}
    	\begin{beamercolorbox}[wd=10mm,ht=4mm,center]{color_Logo-Platzhalter}
   	 		\fontsize{3.5}{5} \selectfont  Platzhalter- \newline Logo  \vspace*{0.3mm}
    	\end{beamercolorbox} 
    \end{minipage}
    \begin{minipage}{0.6\textwidth}
    	\fontsize{4.5}{5.5} \selectfont \color{red} 
    	Für mehr als drei Logos (oder wenn der Platz unten rechts nicht ausreicht) 
    \end{minipage}
{\color{red} \hrule}


\vspace*{1mm}
\begin{flushright}    
    \begin{minipage}{58mm}
    	\begin{flushright}    
    		\begin{minipage}{0.5\textwidth}
				\begin{beamercolorbox}[wd=14mm,ht=4mm,center]{color_Logo-Platzhalter}
    				\fontsize{3.5}{5} \selectfont  Platzhalter- \newline Logo  \vspace*{0.3mm}
    			\end{beamercolorbox} 
    		\end{minipage}
    		\begin{minipage}{0.5\textwidth}
				\begin{beamercolorbox}[wd=26mm,ht=4mm,center]{color_Logo-Platzhalter}
    				\fontsize{3.5}{5} \selectfont  Platzhalter- \newline Logo  \vspace*{0.3mm}
    			\end{beamercolorbox} 
    		\end{minipage}
   			\begin{minipage}{0.5\textwidth}
				\begin{beamercolorbox}[wd=14mm,ht=4mm,center]{color_Logo-Platzhalter}
    				\fontsize{3.5}{5} \selectfont  Platzhalter- \newline Logo  \vspace*{0.3mm}
   				\end{beamercolorbox} 
    		\end{minipage}
			{\color{red} \hrule}
    	\end{flushright}
    \end{minipage}
\end{flushright}

\end{overlayarea}
\end{frame}
%---------------------------------------------------------------------------------
%-----Seite 3---------------------------------------------------------------------
%---------------------------------------------------------------------------------  
\begin{frame}

%\vspace*{5mm}

\textcolor{uolblue}{\fontsize{16}{17}\selectfont{}Kapitelüberschrift}
\bigskip

	\begin{tikzpicture}
    	\fill[gr] (2,0) rectangle (15,5);
	\end{tikzpicture}
 

\end{frame}
%---------------------------------------------------------------------------------
%-----Seite 4---------------------------------------------------------------------
%---------------------------------------------------------------------------------
\begin{frame}{\vspace*{10mm}Überschrift}
\vspace*{-10mm}
\begin{itemize}
\item Aufzählung Standard
	\begin{itemize}
	\item
	Aufzählung Ebene 2, zweizeilig\\
	Aufzählung Ebene 2
		\begin{itemize}
		\item Aufzählung Ebene 3
		\item ...
		\end{itemize}
	\item ...
	\end{itemize}
     	\item Aufzählung Standard
     	\item \fontsize{11}{16}\selectfont Aufzählung vergrößert\\
     	 zweizeilig
\end{itemize}
   
\end{frame}
%---------------------------------------------------------------------------------
%-----Seite 5---------------------------------------------------------------------
%---------------------------------------------------------------------------------
\begin{frame}{\vspace*{10mm}Überschrift}

\lipsum[1][1-6]\\

\begin{center}
\begin{tikzpicture}
    \draw[very thick,blue2] (15,0) -- (16.1,0.4);
    \draw[very thick,blue2] (15,0) -- (13.9,0.4);
  \end{tikzpicture}
  
{\color{blue2} \fontsize{16}{16} \selectfont%
Dies ist der Blindtext einer\\[0.2ex]
farbigen Hervorhebung, wie zum Beispiel\\[0.2ex]
eine Zusammenfassung.}
\end{center}
\end{frame}
%---------------------------------------------------------------------------------
%-----Seite 6---------------------------------------------------------------------
%---------------------------------------------------------------------------------
\begin{frame}{\vspace*{10mm}Überschrift}

\begin{multicols}{2} 

Fließtext, 2-Spaltig. \lipsum[2][1-4]

\medskip

\textbf{Eine Zwischenüberschrift}\\

\lipsum[2][1-4] \textbf{Eine Hervorhebung.} \lipsum[2][5-11]
\end{multicols}
\end{frame}
%---------------------------------------------------------------------------------
%-----Seite 7---------------------------------------------------------------------
%---------------------------------------------------------------------------------
\begin{frame}{\vspace*{10mm}Überschrift}
\begin{multicols}{2} 
\lipsum[2][1-5]

\medskip

\lipsum[2][6-10]
\begin{center}\vspace{1cm}
%
%-Für echte Farben in Ihrem Bild soll 'decodearray={1 0 1 0 1 0}' entfernt werden-
\includegraphics[width=0.8\linewidth,decodearray={1 0 1 0 1 0}]{UOL_Logo_Negativ}
%
\end{center}\vspace{1cm}
\end{multicols}
\end{frame}
%---------------------------------------------------------------------------------
%-----Seite 8---------------------------------------------------------------------
%---------------------------------------------------------------------------------
\begin{frame}

\renewcommand{\arraystretch}{1.5}
\begin{tabular}[h]{p{8cm}cr}
\arrayrulecolor{blue2}\hline
\textbf{Überschrift einer Tabelle}\\
\hline
Zeile einer Tabelle & 00 AB & Lorem Ipsum \\
Zeile einer Tabelle & 00 AB & Lorem Ipsum \\
Zeile einer Tabelle & 00 AB & Lorem Ipsum \\
Zeile einer Tabelle & 00 AB & Lorem Ipsum \\
\hline
\end{tabular}

\bigskip

\begin{tabular}[h]{lcr}
\arrayrulecolor{blue2}\hline
\textbf{Überschrift einer Tabelle}\\
\hline
Zeile einer Tabelle && 00 AB \\
Zeile einer Tabelle && 00 AB \\
Zeile einer Tabelle && 00 AB \\
Zeile einer Tabelle && 00 AB \\
Zeile einer Tabelle && 00 AB \\
\hline
\end{tabular}

\end{frame}
%---------------------------------------------------------------------------------
%-----Seite 9---------------------------------------------------------------------
%---------------------------------------------------------------------------------
\begin{frame}


\mbox{\colorbox{blue2}{
\begin{minipage}{0.25\textwidth}
	\begin{center}
	\color{white} 
	Beispiel für\\
	zentrierten Blindtext\\ 
	auf einem\\
	farbigen Infokasten
	\end{center}
\end{minipage}
}
\quad
\fcolorbox{uolblue}{white}{
\begin{minipage}{0.25\textwidth}
	\color{uolblue} 
	Beispiel für links-\\
	bündigen Blindtext\\ 
	auf einem Outline-\\
	Infokasten
\end{minipage}
}
\quad
\begin{minipage}{0.15\textwidth}
\begin{tikzpicture}[yscale=0.3]
    \draw [line width=1mm,    blue2] (0,-1) -- (2,-1);
    \draw [line width=0.75mm, blue2] (0,-2) -- (2,-2);
    \draw [line width=0.5mm,  blue2] (0,-3) -- (2,-3);
    \draw [line width=0.3mm,  blue2] (0,-4) -- (2,-4);
    \draw [line width=0.2mm,  blue2] (0,-5) -- (2,-5);
    \draw [line width=0.1mm,  blue2] (0,-6) -- (2,-6);
  \end{tikzpicture}
\end{minipage}
\quad
\begin{minipage}{0.15\textwidth}
\begin{tikzpicture}[yscale=0.3]
    \draw [line width=1mm,    uolblue] (0,-1) -- (2,-1);
    \draw [line width=0.75mm, uolblue] (0,-2) -- (2,-2);
    \draw [line width=0.5mm,  uolblue] (0,-3) -- (2,-3);
    \draw [line width=0.3mm,  uolblue] (0,-4) -- (2,-4);
    \draw [line width=0.2mm,  uolblue] (0,-5) -- (2,-5);
    \draw [line width=0.1mm,  uolblue] (0,-6) -- (2,-6);
  \end{tikzpicture}
\end{minipage}}


\begin{flushleft}
\smartdiagramset{uniform sequence color=true, 
sequence item uniform color=blue2,
sequence item border color=white,
sequence item font size=\footnotesize, 
sequence item text color=white,
sequence item height=1.5cm,
sequence item width=1.8cm
}
\smartdiagram[sequence diagram]{1,2,3,4} 
\end{flushleft}


Keine Verläufe
 
Keine farbigen Schriften auf farbigen Flächen
 
Keine Schrift auf Bildern
 
Grundsätzliches Farbklima: Blau

\end{frame}
%---------------------------------------------------------------------------------
%-----Seite 10---------------------------------------------------------------------
%---------------------------------------------------------------------------------
\begin{frame}{\vspace*{10mm}Farben}
\vspace*{-10mm}
\begin{multicols}{2} 
\textbf{Primärfarben}

\bigskip

UOL-Blau

\begin{tikzpicture}
    \fill[uolblue]
  (2,0) rectangle (7.9,1.5);
  \end{tikzpicture}
  
\medskip
  
Farbgruppe Blau

\begin{tikzpicture}
    \fill[blue1]
  (2,0) rectangle (3.4,1.5);
    \fill[blue2]
  (3.5,0) rectangle (4.9,1.5);
    \fill[blue3]
  (5,0) rectangle (6.4,1.5);
    \fill[blue4]
  (6.5,0) rectangle (7.9,1.5);
  \end{tikzpicture}
  
\textbf{Sekundärfarben}

\bigskip

Farbgruppe Grün

\begin{tikzpicture}
     \fill[green1]
  (2,0) rectangle (3.4,1.5);
    \fill[green2]
  (3.5,0) rectangle (4.9,1.5);
    \fill[green3]
  (5,0) rectangle (6.4,1.5);
    \fill[green4]
  (6.5,0) rectangle (7.9,1.5);
  \end{tikzpicture}

\medskip
  
Farbgruppe Orange

\begin{tikzpicture}
    \fill[orange1]
  (2,0) rectangle (3.4,1.5);
    \fill[orange2]
  (3.5,0) rectangle (4.9,1.5);
    \fill[orange3]
  (5,0) rectangle (6.4,1.5);
    \fill[orange4]
  (6.5,0) rectangle (7.9,1.5);
  \end{tikzpicture}
\end{multicols}
\end{frame}
 \end{document}